% ==============================================================================
% banq-bundle.tex
% Defines the \ifbundle conditional that distinguishes a standalone child build
% from the unified bundle build, and the \xref macro that produces a real
% internal hyperlink in bundle mode and a citation-with-text in standalone mode.
%
% Standalone child papers leave \bundlefalse (the default).
% The master bundle root sets \bundletrue before \begin{document}.
% ==============================================================================

\newif\ifbundle
\bundlefalse

% \xref{<bibkey>}{<label>}{<text>}
% - In bundle mode: produces a clickable in-document hyperlink to the label.
% - In standalone mode: prints the text followed by a citation to the bibkey.
% Note: defined after \bundletrue/\bundlefalse so it picks up the right branch.
% Children must call this from inside \AtBeginDocument to ensure \ifbundle is
% finalized; alternatively the macro can test \ifbundle at call site.
\newcommand{\xref}[3]{%
  \ifbundle
    \hyperref[#2]{#3}%
  \else
    #3~\cite{#1}%
  \fi
}

% \xrefshort{<bibkey>}{<label>}{<text>}
% Same as \xref but suppresses the citation entirely in standalone mode (used
% when the surrounding prose already cites the source).
\newcommand{\xrefshort}[3]{%
  \ifbundle
    \hyperref[#2]{#3}%
  \else
    #3%
  \fi
}

% \banqEnableBundleNumbering
% Resets the section counter at every \part and prefixes \thesection with the
% part's Roman numeral, so headings read "I.1 Introduction", "II.1 Introduction",
% etc. Called explicitly from the master root after \bundletrue (see
% 000.banq-all/banq.tex). Not called by standalone children.
%
% Parts that contain MULTIPLE sub-papers must additionally call \banqSubpart
% (defined below) before each child \subfile, which appends a per-paper letter
% suffix (a, b, ...) to the part numeral and resets the section counter so each
% sub-paper restarts from 1. Without this, sub-papers within the same part
% share a single section sequence, which causes the TOC to conflate identically
% titled sections (e.g. two "II.1 Introduction" entries for banq-lck/banq-log).
%
% The outer \makeatletter/\makeatother brackets here are critical: the macro
% body contains \@addtoreset (an @ control sequence) which must be parseable
% as a single token at \newcommand-definition time, not at call time.
\newcommand{\banqSubpartLetter}{}
% Boolean tracking whether we are currently inside a sub-paper of a shared part
% (i.e. \banqSubpart has been called since the last \part). Used by \banqPaper
% to decide whether to demote section bookmark levels. \ifx-comparing
% \banqSubpartLetter against \@empty doesn't work because \newcommand produces
% \long macros while \@empty is short — \ifx always returns false.
\newif\ifbanqSharedPart
\makeatletter
\newcommand{\banqEnableBundleNumbering}{%
  \@addtoreset{section}{part}%
  \renewcommand{\thesection}{\Roman{part}\banqSubpartLetter.\,\arabic{section}}%
  % Make hyperref anchor names unique across sub-papers that share a \part.
  % Without this, banq-lck and banq-log both emit section.2.1, section.2.2, …
  % which corrupts the PDF outline tree (external viewers truncate at Part II).
  \renewcommand{\theHsection}{\Roman{part}\banqSubpartLetter.\arabic{section}}%
  \renewcommand{\theHsubsection}{\theHsection.\arabic{subsection}}%
  \renewcommand{\theHsubsubsection}{\theHsubsection.\arabic{subsubsection}}%
  \renewcommand{\theHparagraph}{\theHsubsubsection.\arabic{paragraph}}%
  % Hook \part so the subpart-letter and shared-part flag both reset when a new
  % part begins; otherwise state set in part N would leak into part N+1.
  \let\banqOriginalPart\part
  \renewcommand{\part}{%
    \renewcommand{\banqSubpartLetter}{}%
    \banqSharedPartfalse
    \banqOriginalPart
  }%
  % The "I.\,10" prefix is wider than the default \numberline box. Widen the
  % TOC number-box widths so two-digit section numbers don't overlap titles.
  \setlength{\cftsecnumwidth}{3.5em}%
  \setlength{\cftsubsecnumwidth}{4.2em}%
  \setlength{\cftsubsubsecnumwidth}{5.2em}%
}
\makeatother

% \banqSubpart{<letter>}
% Marks the start of a sub-paper inside a part that contains multiple papers.
% Sets the per-paper letter suffix appended to the part numeral in \thesection
% (so headings read "IIa.1 Introduction", "IIb.1 Introduction", etc.) and
% resets the section counter so the sub-paper's own numbering restarts at 1.
% Pass an empty argument (\banqSubpart{}) to clear the suffix mid-part if
% needed; the suffix also auto-clears at every \part boundary (see hook above).
\newcommand{\banqSubpart}[1]{%
  \renewcommand{\banqSubpartLetter}{#1}%
  \setcounter{section}{0}%
  \banqSharedParttrue
}

% \banq@RestoreLevels / \banq@DemoteLevels
% Adjust hyperref's per-sectioning bookmark levels. When two sub-papers share a
% \part (banq-lck "IIa" and banq-log "IIb" inside Part~II), each paper's
% sections should nest under its own divider bookmark in the PDF outline,
% rather than appearing as siblings of it under the part. Demoting
% \toclevel@section from 1 to 2 (and the deeper levels accordingly) achieves
% that nesting; \banqPaper calls Restore before emitting its own divider so the
% divider always lands at level 1, then calls Demote afterwards (only when in a
% shared part) so subsequent sections sit at level 2 = child of the divider.
% Affects only the PDF outline depths; the printed TOC layout is unchanged.
\makeatletter
\newcommand{\banq@RestoreLevels}{%
  \renewcommand{\toclevel@section}{1}%
  \renewcommand{\toclevel@subsection}{2}%
  \renewcommand{\toclevel@subsubsection}{3}%
  \renewcommand{\toclevel@paragraph}{4}%
}
\newcommand{\banq@DemoteLevels}{%
  \renewcommand{\toclevel@section}{2}%
  \renewcommand{\toclevel@subsection}{3}%
  \renewcommand{\toclevel@subsubsection}{4}%
  \renewcommand{\toclevel@paragraph}{5}%
}
\makeatother

% Storage for the current part's twemoji name. \banqPartWithEmoji sets it
% (overwritten at each part boundary). \banqPaper reads it without clearing,
% so multiple sub-papers within the same part (e.g. banq-lck and banq-log
% in Part~II) share the same emoji on their respective cover pages.
\makeatletter
\def\banq@partemoji{}
\newif\ifbanq@hasemoji
\banq@hasemojifalse
% \banqSetSubpartEmoji{<twemoji-name>}
% Override the part's default emoji for the next sub-paper cover. Used inside
% parts that contain multiple papers when each sub-paper deserves its own
% cover icon (e.g. banq-lck takes "lock" while banq-log keeps the part-level
% "gear"). Encapsulates the @-internal \gdef so the master document body need
% not enter \makeatletter scope.
\newcommand{\banqSetSubpartEmoji}[1]{\gdef\banq@partemoji{#1}}
\makeatother

% \banqForceRecto
% Page-break helper that flushes outstanding output and ensures the next
% content-bearing page lands on an odd (recto) page in duplex bindings.
% The article class is loaded without `twoside` so \cleardoublepage degenerates
% to \clearpage and will not pad; we check the page counter explicitly and
% insert a blank filler when the upcoming page is even. \value{page} reflects
% the next page number after \clearpage, before any content has been emitted
% on it. Call sites: front-matter (TOC start, body start) and \banqPaper.
\newcommand{\banqForceRecto}{%
  \clearpage
  \ifodd\value{page}\else\hbox{}\thispagestyle{empty}\newpage\fi
}

% \banqEnsureEvenPageCount
% Final-pass parity helper: guarantees the document ends with an even total
% page count, so duplex-bound copies have a back cover on the verso side and
% any subsequent insert (e.g. a colophon, a bound-in errata sheet) starts on
% a recto. After \clearpage, \value{page} holds the *next* page number, so
% the current page count equals \value{page}-1; the count is even iff
% \value{page} is odd, which is exactly the invariant \banqForceRecto already
% enforces. The two operations are semantically distinct (recto-start vs.
% even-total) but compute identical pad logic, so we alias them. Call once,
% immediately before \end{document}, in the bundle root.
\let\banqEnsureEvenPageCount\banqForceRecto

% \banqPartWithEmoji{<title>}{<twemoji-name>}
% Bumps the part counter, stashes the title for the part-level TOC entry,
% and stashes the emoji for the cover page(s) that follow (see \banqPaper).
% No physical divider page is emitted — the cover that immediately follows
% is the visible boundary between parts. The bundle's section-numbering
% scheme (I.1, IIa.1, IIb.1, III.1, ...) continues to work because
% \refstepcounter drives the same counter \part would have, and
% \@addtoreset (set up by \banqEnableBundleNumbering) still resets section
% at the part boundary.
%
% The TOC entry itself is written by \banqPaper at the start of the cover
% page rather than here, so the deferred page-number capture lands on the
% cover (the visible part boundary) rather than on whatever page the build
% happens to be on between subfiles.
\makeatletter
\def\banq@parttitle{}
\newif\ifbanq@parttocpending
\banq@parttocpendingfalse
\newcommand{\banqPartWithEmoji}[2]{%
  \refstepcounter{part}%
  \renewcommand{\banqSubpartLetter}{}%
  \banqSharedPartfalse
  \gdef\banq@partemoji{#2}%
  \global\banq@hasemojitrue
  \gdef\banq@parttitle{#1}%
  \global\banq@parttocpendingtrue
}
\makeatother

% \banqAbstract{<abstract text>}{<keywords>}
% Single source of truth for each child's abstract + keywords list, with
% mode-aware rendering:
%   - Standalone: emits a normal \begin{abstract}...\end{abstract} block
%     followed by a "Keywords:" line (the previous \ifbundle\else ... \fi
%     boilerplate each child used to inline).
%   - Bundle: stashes the content into module-level token lists and sets
%     \ifbanq@hasabstract; \banqPaper then renders it on the verso page that
%     follows the cover (the page that was previously left blank).
% Defined as \long so abstracts containing paragraph breaks are accepted,
% though in practice all six abstracts are single-paragraph.
\makeatletter
\def\banq@abstracttext{}
\def\banq@keywordstext{}
\newif\ifbanq@hasabstract
\banq@hasabstractfalse
\long\def\banqAbstract#1#2{%
  \ifbundle
    \long\gdef\banq@abstracttext{#1}%
    \long\gdef\banq@keywordstext{#2}%
    \global\banq@hasabstracttrue
  \else
    \begin{abstract}#1\end{abstract}\par\medskip\noindent
    \textbf{Keywords:} #2\par
  \fi
}
\makeatother

% \banqPaper{<title>}{<subtitle>}
% Paper-level divider page printed at the top of each child's body when in
% bundle mode. Sits below the \part page in the visual hierarchy and provides
% a clear transition between consecutive papers within the same part (e.g.
% banq-lck -> banq-log inside Part~II). A no-op in standalone builds.
%
% Wrapped in \makeatletter so the macro body can reference \banq@RestoreLevels,
% \banq@DemoteLevels, and \@empty as single tokens rather than \banq + @… .
\makeatletter
\newcommand{\banqPaper}[2]{%
  \ifbundle
    % Force the cover page to land on an odd (recto) page; see \banqForceRecto.
    \banqForceRecto
    \onecolumn
    \thispagestyle{empty}
    % If a \banqPartWithEmoji call is pending, write its TOC + PDF-outline
    % entry now so the captured page number is the cover page (the visible
    % part boundary), not whatever blank pad the build happened to land on.
    \ifbanq@parttocpending
      \phantomsection
      \addcontentsline{toc}{part}{\thepart\hspace{1em}\banq@parttitle}%
      \global\banq@parttocpendingfalse
    \fi
    \vspace*{6em}
    \begin{center}
      {\Large\color{codegray} \textsc{XPower Banq: Part \Roman{part}\banqSubpartLetter}}\\[0.6em]
      \rule{0.4\textwidth}{0.4pt}\\[1.2em]
      {\fontsize{22}{26}\selectfont\bfseries #1}\\[0.8em]
      {\large\color{codegray} #2}
    \end{center}
    \vspace*{\fill}
    \ifbanq@hasemoji
      \begin{center}
        \twemoji[height=8cm]{\banq@partemoji}
      \end{center}
    \fi
    \vspace*{\fill}
    \clearpage
    % Verso page after the cover. If the child supplied an abstract via
    % \banqAbstract, render it here (still in onecolumn, unstyled body width
    % so the prose breathes). Otherwise leave the page blank — used for the
    % "References & Glossary" cover, which has no per-paper abstract since
    % the bundle's front-matter abstract already summarises the whole volume.
    \thispagestyle{empty}
    \ifbanq@hasabstract
      \vspace*{\fill}
      \begin{center}{\large\bfseries Abstract}\end{center}
      \medskip
      \begingroup
        \small
        \noindent\banq@abstracttext\par
        \medskip
        \noindent\textbf{Keywords:} \banq@keywordstext\par
      \endgroup
      \vspace*{\fill}
      \global\banq@hasabstractfalse
    \else
      \hbox{}%
    \fi
    \newpage
    \twocolumn
    % Register the divider in the TOC and PDF outline as a child of the part.
    % Restore section bookmark levels first so this divider always lands at
    % level 1 (regardless of any prior demotion left over from a sibling
    % sub-paper), then conditionally demote so subsequent sections nest under
    % this divider when we're inside a shared part.
    \banq@RestoreLevels
    \phantomsection
    \addcontentsline{toc}{section}{\protect\numberline{\Roman{part}\banqSubpartLetter.}#1}%
    \ifbanqSharedPart
      \banq@DemoteLevels
    \fi
  \fi
}
\makeatother
